\subsection{Hecke 算子与翻译器：\S9--\S12（pp.123--222）}

本部分覆盖原文第 9 至第 12 节（印刷页 pp.123--222），主题是 't~Hooft 算子、Hecke 修改以及相应 A-膜在四维规范理论中的构造。问答只取原文用陈述句明确给出定义、结论或公式的地方（A 级材料）；凡原文只是暗示、没有明说的推论一律不录。每条答案之后给出印刷页锚点与英文短引（不超过 120 字符），以便逐条回查原文。

\QAq{\S9 为什么值得单独展开？}
\QAa{原文在本节要证明，$\Psi=0$ 处的超对称 't~Hooft 算子恰好对应几何 Langlands 的 Hecke 算子；$\Psi=0$ 由取 $t=1$、$\theta=0$ 达到。此时四维 TQFT 在二维化后成为以 $\MH$ 为目标、复结构 $K$ 的 A-模型；A-模型没有瞬子修正，超对称位形都时间无关。}
\QAs{src p.123}{``'t~Hooft operators correspond to the Hecke operators of the geometric Langlands program.''}

\QAq{'t~Hooft 算子是普通意义下的算子吗？}
\QAa{不是。原文强调静态 't~Hooft 算子支撑在 $\RR\times p$ 上，属于“线算子”，它们必须被并入物理的、超对称的态空间的定义之中，而不能当作通常 Hilbert 空间上的算子。这一定位正是后文把 't~Hooft 算子与 Hecke 修改对应起来的前提。}
\QAs{src p.123}{``they must be incorporated in the definition of the space of physical, supersymmetric states.''}

\QAq{$U(1)$ 情形 't~Hooft 算子怎样作用？}
\QAa{算子 $T(m;p_0)$ 把线丛 $L$ 变为 $L\otimes\mathcal O(p_0)^m$，即按 $\mathcal O(p_0)$ 的 $m$ 次幂作扭化。原文指出这与 $U(1)$ 的 Hecke 算子的标准定义一致；同一 $p_0$ 处的算子构成交换群 $T(m)=T(1)^m$，不同点的算子也彼此交换。}
\QAs{src p.131}{``This result agrees with the standard definition of the Hecke operators for U(1).''}

\QAq{$U(N)$ 情形的奇点与 Hecke 修改如何对应？}
\QAa{'t~Hooft 算子由同态 $\rho:U(1)\to U(N)$（等价于整数权 $m_1\ge\cdots\ge m_N$）的共轭类分类；在 $p_0$ 附近 $E$ 分解为 $L_1\oplus\cdots\oplus L_N$，修改把 $L_i$ 换成 $L_i\otimes\mathcal O(p_0)^{m_i}$。原文说这与 $U(N)$ 的 Hecke 算子的数学描述完全一致，并称之为 't~Hooft 算子与 Hecke 算子之间的基本联系。}
\QAs{src p.132}{``This coincides with the usual mathematical description of the action of the Hecke operators for U(N).''}

\QAq{原文所说的 Hecke 修改究竟指什么？}
\QAa{一个 Hecke 修改是带指定同构的丛对：$E_+$ 在 $p_0$ 之外与 $E_-$ 自然同构，而这一同构一般不能延拓到 $p_0$。局部地看 $E_+=\oplus_i\,(L_i\otimes\mathcal O(p_0)^{m_i})$；不同的局部分解可能给出同一个 $E_+$，原文正是要研究这些分解商掉等价关系之后的空间。}
\QAs{src p.136}{``any Hecke modification $E_+$ of $E_-$ is naturally isomorphic to $E_-$ away from $z = 0$''}

\QAq{Hecke 修改的空间为何需要紧化？}
\QAa{A-模型的物理态是这些修改空间上的上同调，原文指出紧性使相应上同调群的定义简单直接。$U(2)$ 的权 $(1,0)$ 给出光滑紧的 $\CC\PP^1$；而 $(2,0)$ 的自然紧化必须额外允许 $(1,1)$ 型修改，后者就是 $E_-\otimes\mathcal O(p_0)$。}
\QAs{src p.138}{``Compactness makes the definition of the appropriate cohomology groups much more straightforward''}

\QAq{仿射 Grassmannian 与 Hecke 修改空间是什么关系？}
\QAa{带 $p_0$ 外平凡化的丛对构成商 $GL(N,\mathcal X)/GL(N,\mathcal O)$，即仿射 Grassmannian $\mathrm{Gr}_N$；它同构于 $GL(N,\CC((z)))/GL(N,\CC[[z]])$，也同构于 $U(N)$ 的基于点的环路空间。权 $Lw=(m_1,\dots,m_N)$ 对应的修改空间是点 $z^{-m_j}\delta_{ji}$ 在 $GL(N,\mathcal O)$ 下的轨道，这些轨道（Schubert 胞腔）恰好分层整个 $\mathrm{Gr}_N$。}
\QAs{src p.146}{``The latter space is known as the affine Grassmannian $\mathrm{Gr}_N$ for the group $GL(N)$.''}

\QAq{\S10.1 在三维流形上取什么边界条件？}
\QAa{左端固定 $E_-$（即对 $A$ 取 Dirichlet 条件，$\varphi_0$ 自由），右端令 $\varphi_0=0$、$A$ 满足协变 Neumann 条件 $\iota_nF=0$，并且只商去右端非平凡的规范变换。这样定义的椭圆复形的指标给出修改空间的虚维数；无 't~Hooft 算子时该模空间退化为一点。}
\QAs{src p.150}{``we want to allow all possible Hecke modifications $E_+$ that may arise from $E_-$ by solving the Bogomolny equations''}

\QAq{单极子冒泡是什么？}
\QAa{等变瞬子在三维图像中可以把一个量子化单位的电荷集中到一点并使其消失，此即单极子冒泡。原文指出它只能发生在存在 't~Hooft 算子的情形；ADHM 分析表明冒泡把权 $(m+s,m-1)$ 逐步降为 $(m+s-1,m)$，这正是 \S9.2 中 associated 权所描述的退化。}
\QAs{src p.156}{``Monopole bubbling can occur only in the presence of an 't~Hooft operator''}

\QAq{\S10.3 的 Kähler 结构与 moment map 起什么作用？}
\QAa{原文在场空间上引入复结构 $I$（坐标为 $A_z$ 与 $B=A_y-i\varphi_0$）、辛形式与 Kähler 度规；Bogomolny 方程于是写成全纯方程加 moment map 方程 $\mu=0$。$\mu=0$ 还要求边界上 $\Tr\epsilon\varphi_0=0$，这恰好由左端固定规范、右端令 $\varphi_0=0$ 实现，所以原文称这两种边界条件是最简单自然的选择。又因规范变换在左端平凡，$G_\CC$ 自由作用，辛商与复商相同：$Z=Z'$。}
\QAs{src p.161}{``they are the simplest ones that allow us to interpret the Bogomolny equations in terms of a holomorphic equation''}

\QAq{\S10.4 中两个 't~Hooft 算子的 OPE 是什么？}
\QAa{当 $p'\to p$ 时，乘积 $T(Lw,p)\cdot T(Lw',p')$ 展开为 $T(Lw+Lw',p')$ 再加上所有 associated 权 $Lw_i$ 的项。原文说这是 't~Hooft 算符乘积展开的一部分；取 $B,B'$ 为复结构 $I$ 中的复 Lagrangian 子流形（即 $(B,A,A)$ 膜）时，超对称构型空间正是 $Z(LR_1,\dots,LR_k)$，物理态取其上同调。}
\QAs{src p.165}{``this statement is part of the operator product expansion of the operators $T(Lw,p)$ and $T(Lw',p')$''}

\QAq{为什么必须放弃时间反演不变性？}
\QAa{时间反演不变性只适用于所取膜使解满足 $\phi=0$ 的情形；几何 Langlands 最重要的 A-膜（Hitchin 纤维与 c.c. 膜）并不限于 $\phi=0$。因此研究 't~Hooft 算子对它们的作用时必须放弃该假设，改用扩展 Bogomolny 方程。其解沿 $y$ 方向只差规范变换，直到穿过 't~Hooft 算子处发生 Higgs 丛的 Hecke 修改。}
\QAs{src p.179}{``requiring time-reversal invariance puts a severe restriction on the allowed boundary conditions''}

\QAq{Higgs 丛的 Hecke 修改为何是有限的？}
\QAa{全纯丛可以在一点处修改而保持全纯性，但全纯场 $\varphi$ 不能：$\varphi$ 在 $z_0$ 之外给定后，至多只有唯一的全纯延拓。原文算出所需条件数为 $\sum_{i<j}(m_i-m_j)$，恰好等于 (9.25) 给出的修改空间维数，因此 Higgs 丛的 Hecke 修改通常是有限集；当 $\varphi(z_0)$ 正则半单时确实有限，此时 $\varphi$ 在 $Y(Lw,z_0)$ 上定义向量场，其零点就是可能的修改。}
\QAs{src p.180}{``including $\varphi$ generically eliminates some parameters''}

\QAq{典范余各向同性膜是什么？}
\QAa{原文先定义余各向同性子流形：它局部由 Poisson 交换的一阶约束函数定义，维数至少为一半。取支撑为整个 $\MH$、Chan-Paton 为秩一线丛时，A-膜条件 $(\omega^{-1}F)^2=-1$ 给出的复结构可积，且 $\omega-iF$ 必须是全纯辛形式，故这类膜只当 $\omega$ 是全纯辛形式的实部时才可能存在。$\MH$ 上取 $F=(\Im\tau)\omega_J$（此时 $\hat N=I$）所定义的膜称为典范余各向同性膜（c.c. 膜）；它同时是 $(A,B,A)$ 膜，即在复结构 $I$、$K$ 中是 A-膜、在 $J$ 中是 B-膜。}
\QAs{src p.186}{``We will call the A-brane on $M_H$ with $F=(\Im\tau)\omega_J$ the canonical coisotropic brane''}

\QAq{c.c. 膜对应什么 D-模？}
\QAa{c.c. 膜的存在意味着复结构 $K$ 中每个 A-膜都是 $K_M^{1/2}$ 上微分算子 sheaf 的模；任意两个膜的弦构成相应代数上的模，这是开闭 TFT 的公理。对 $\varphi=0$ 的 $(B,A,A)$ 膜，BRST 算符 $Q$ 就是 $\bar\partial$，物理弦 sheaf 是 $K_M^{1/2}$ 的全纯截面 sheaf。type F 膜（Hitchin 纤维加平坦线丛）则给出 $\MH$ 上的 D-模 sheaf。}
\QAs{src p.198}{``it is a magnetic eigenbrane because it is the S-dual of a zerobrane; and it gives rise to a sheaf of D-modules''}

\QAq{一般 $\Psi$ 下还有 c.c. 膜吗？}
\QAa{有。原文证明对任意 $\Psi$ 存在类似物 $B^\Psi_{\mathrm{c.c.}}$，其开弦代数是 $K_M^{1/2}\otimes L^\Psi$ 上的全纯微分算子 sheaf。作为代价，当 $\theta\ne2\pi n$ 时 Hitchin 纤维不再是合法的 A-膜，因为 $B|_Y$ 的周期不是 $2\pi$ 的整数倍；点也永远不可能是 A-膜。}
\QAs{src p.202}{``We will see that an analogue of the c.c. brane exists for all $\Psi$.''}

\QAq{扭化微分算子与普通微分算子如何联系？}
\QAa{对整数 $\Psi=n$，$(B^n_{\mathrm{c.c.}},B_n)$ 的弦态代数恰是 $K_M^{1/2}\otimes L^n$ 上的微分算子；再由扰动论证 $f(\Psi)$ 是线性函数，可推出任意 $\Psi$ 时代数为 $L^{\Psi-h}$ 上的全纯微分算子（$h$ 为对偶 Coxeter 数）。取 $\Psi=h$ 即得作用在函数上的普通微分算子，于是 S-对偶加上 $T^h$ 把零膜映为普通的 D-模。}
\QAs{src p.206}{``if we let $\Psi = h$, the algebra of c.c. strings is the sheaf of ordinary differential operators''}

\QAq{四类膜如何从规范理论得到？}
\QAa{原文列出四类：零膜为 $(B,B,B)$，其镜像是 Hitchin 纤维、为 $(B,A,A)$，c.c. 膜为 $(A,B,A)$，其镜像亦然。四类膜都保两个拓扑超荷（连同其伴随超荷共四个超对称，在规范理论中称 $1/4$ BPS 膜，在 sigma-模型语言中称 $1/2$ BPS 膜）。原文强调这些重要膜都可以在四维规范理论中定义为边界条件，尽管提升并不唯一。}
\QAs{src p.207}{``We will focus on these four kinds of brane and describe the corresponding boundary conditions in the gauge theory.''}

\QAq{零膜 $(B,B,B)$ 的边界条件如何确定？}
\QAa{取 $T=\Gamma^{04}$、$T\lambda|=\lambda|$，得到 $F_{0I}+F_{4I}=0$；再令 $A_0+A_4=0$、其余 $A_I$ 取 Dirichlet 条件，并要求存在超对称态，就要求 Hitchin 方程成立。边界值于是在 $\MH$ 中定义一点 $p$，这正是支撑在 $p$ 的零膜；对不可约 Higgs 丛，这一边界条件还被唯一确定。}
\QAs{src p.213}{``The boundary conditions we have obtained so far correspond to a zerobrane on $M_H$ supported at the point $p$.''}

\QAq{$(B,A,A)$ 的边界条件给出什么？}
\QAa{对 $(B,B,B)$ 作 S-对偶，取 $T'=\Gamma^1\Gamma^2\Gamma^3\Gamma^4$、$T'\lambda|=-\lambda|$：$\varphi_z$ 的特征多项式固定，而决定 $E$ 全纯结构的 $A_z$ 可以涨落。再要求 $\Gamma^{2367}\epsilon=\epsilon$，即得 $D_z\varphi_z=0$（复 Hitchin 方程），膜支撑在 Hitchin 纤维上，是典型的 $(B,A,A)$ 膜。}
\QAs{src p.216}{``the field $A_z$ which determines the holomorphic structure of $E$ is allowed to fluctuate''}

\QAq{$(A,B,A)$ 的边界条件给出什么？}
\QAa{取 $T'=-\Gamma^{15}\Gamma^{2-7}\Gamma^{3+6}$ 并令 $T'\lambda|=\lambda|$，得 $A_z+i\varphi_z$ 固定、$A_{\bar z}$ 自由，从而在 $J$ 中定义复的、中间维数的、$\Omega_J$ 迷向的子流形 $W'$。$W'$ 正是固定全纯结构的平坦 $G_\CC$ 丛，故其上的膜是 $(A,B,A)$ 膜；S-对偶后得到 target 空间为整个 $\MH$ 的 space-filling $(A,B,A)$ 膜。}
\QAs{src p.218}{``So the brane supported on $W'$, with a trivial Chan-Paton bundle, is an (A, B, A)-brane, as expected.''}

\QAq{$(A,B,A)$ 的规范场是否就是 c.c. 膜的数据？}
\QAa{是。对 S-对偶得到的 space-filling 膜，把低能边界条件与混合 Dirichlet--Neumann 条件比较，可定出靶空间规范场 $f=\omega_J$。原文指出这正是 c.c. 膜的靶空间规范场，因此该边界条件所构造的膜就是典范余各向同性膜。}
\QAs{src p.220}{``This is the target-space gauge field for the canonical coisotropic (A, B, A)-brane''}
